\documentclass[11pt]{article}
\usepackage[margin=1in]{geometry}
\usepackage{amsmath,amssymb,amsfonts,mathtools}
\usepackage{array,booktabs,enumitem,graphicx,xcolor,hyperref}
\usepackage[T1]{fontenc}
\usepackage[utf8]{inputenc}
\title{Collective Field Representation of Nonrelativistic Fermions in (1+1)
  Dimensions}
\author{Source-backed arXiv sample}
\date{}
\begin{document}
\maketitle
\begin{abstract}
A collective field formalism for nonrelativistic fermions in (1+1) dimensions is presented. Applications to the D=1 hermitian matrix model and the system of one-dimensional fermions in the presence of a weak electromagnetic field are discussed.
\end{abstract}
\section{Source Metadata}
This sample is based on arXiv record \texttt{hep-th/9109013}. Authors: Dimitra Karabali.
\section{Extracted Prose 1}
A collective field formalism for nonrelativistic fermions in (1+1) dimensions is presented. The quantum mechanical fermionic problem is bosonized and converted to a second quantized Schr\textbackslash{}"odinger field theory. A formulation in terms of current and density variables gives rise to the collective field representation. Applications to the math expression hermitian matrix model and the system of one-dimensional fermions in the presence of a weak electromagnetic field are discussed. =0
\section{Extracted Prose 2}
=12pt A collective field formalism for nonrelativistic fermions in (1+1) dimensions is presented. The quantum mechanical fermionic problem is bosonized and converted to a second quantized Schr\textbackslash{}"odinger field theory. A formulation in terms of current and density variables gives rise to the collective field representation. Applications to the math expression hermitian matrix model and the system of one-dimensional fermions in the presence of a weak electromagnetic field are discussed.
\end{document}
